% Auto-generated from Chapter-47-Amara-Keeper.md
% Source: C:\Users\PCGAME\Desktop\story&universes\chain-weapon-story\finished-manuscript\manuscriptbaseforchapters\Chapter-47-Amara-Keeper.md

\chapter{Amara Becomes Keeper}
\renewcommand{\CurrentChapterNum}{47}
\POV{\glyphAmara}{Amara}



The ink was the colour of river mud at dusk—not black, not brown, but the intermediate shade that occurred when the ingredients were mixed by hand from oak gall and iron salt and the gum arabic that Amara sourced from a merchant in the delta markets who did not ask why she needed quantities sufficient for a small scriptorium and who accepted payment in the form of intelligence about trade route disruptions, which was, in the economy of the displaced, more valuable than coin.

She made the ink herself. She had always made it herself. The process was part of the discipline—the grinding of the gall, the measuring of the salt, the slow stirring that released the tannins into the solution and produced a pigment whose colour deepened over hours, the fresh ink pale and watery on the page and the aged ink dark and indelible, the words gaining weight as the chemistry completed itself, the facts becoming permanent through a process that required nothing but patience and time.

The encoded journal lay open on the table. The table was a plank of oak set across two barrel staves in the upper room of a factor's warehouse on the river, and the room smelled of the hemp bales stored below—a green, fibrous scent that mixed with the mineral tang of the ink and the warm leather of the journal's binding and the faint musk of the lamp oil that burned in the clay dish at her elbow, the flame low and steady, the wick trimmed to the precise length that produced light sufficient for writing without producing light sufficient for a shadow to be seen from outside the shutter.

Three weeks since the ridgeline. Three weeks since her hands had pressed against the wound that her hands could not close. Three weeks during which the grief had undergone a transformation that she had observed in herself with the clinical attention of a woman trained to observe—the acute phase passing, the rawness calcifying, the open wound becoming scar tissue that was stronger than the surrounding skin but less sensitive, the body's compromise between healing and remembering: you could be repaired or you could feel, but not both, and the body chose repair because survival required it.

She chose both. The griot tradition demanded both.

The bells in her braids chimed as she leaned forward to dip the pen. The small brass bells—seven of them, one for each year of her apprenticeship to the keeper's tradition, strung on threads of indigo cotton that matched the cloth she wore over her shoulders—the bells were the notation system that preceded her literacy, the tradition that said a keeper's first library was her body, and the body's first catalogue was sound. Each bell was tuned differently. Each chime was a register. When she moved, the bells recorded the movement in a chord that was unique to her gait, her posture, the angle of her head, and the chord was a fact, and the fact was a record, and the record was a keeping.

She wrote.